\documentclass{article}%
\usepackage{amsmath}
\usepackage{amsfonts}
\usepackage{amssymb}
\usepackage{graphicx}%
\setcounter{MaxMatrixCols}{30}
%TCIDATA{OutputFilter=latex2.dll}
%TCIDATA{Version=5.50.0.2953}
%TCIDATA{CSTFile=40 LaTeX article.cst}
%TCIDATA{Created=Thursday, May 21, 2026 16:16:11}
%TCIDATA{LastRevised=Thursday, May 21, 2026 17:39:19}
%TCIDATA{<META NAME="GraphicsSave" CONTENT="32">}
%TCIDATA{<META NAME="SaveForMode" CONTENT="1">}
%TCIDATA{BibliographyScheme=Manual}
%TCIDATA{<META NAME="DocumentShell" CONTENT="Standard LaTeX\Blank - Standard LaTeX Article">}
%TCIDATA{ComputeDefs=
%$i_{2}(t)=[\frac{1}{C}\int[i_{1}(t)-i_{2}(t)]dt]\frac{1}{R_{2}}$
%}
%BeginMSIPreambleData
\providecommand{\U}[1]{\protect\rule{.1in}{.1in}}
%EndMSIPreambleData
\newtheorem{theorem}{Theorem}
\newtheorem{acknowledgement}[theorem]{Acknowledgement}
\newtheorem{algorithm}[theorem]{Algorithm}
\newtheorem{axiom}[theorem]{Axiom}
\newtheorem{case}[theorem]{Case}
\newtheorem{claim}[theorem]{Claim}
\newtheorem{conclusion}[theorem]{Conclusion}
\newtheorem{condition}[theorem]{Condition}
\newtheorem{conjecture}[theorem]{Conjecture}
\newtheorem{corollary}[theorem]{Corollary}
\newtheorem{criterion}[theorem]{Criterion}
\newtheorem{definition}[theorem]{Definition}
\newtheorem{example}[theorem]{Example}
\newtheorem{exercise}[theorem]{Exercise}
\newtheorem{lemma}[theorem]{Lemma}
\newtheorem{notation}[theorem]{Notation}
\newtheorem{problem}[theorem]{Problem}
\newtheorem{proposition}[theorem]{Proposition}
\newtheorem{remark}[theorem]{Remark}
\newtheorem{solution}[theorem]{Solution}
\newtheorem{summary}[theorem]{Summary}
\newenvironment{proof}[1][Proof]{\noindent\textbf{#1.} }{\ \rule{0.5em}{0.5em}}
\begin{document}

\begin{thebibliography}{}}                                                                                                %
$Ve(t)=R_{1}i_{1}(t)+L\frac{di_{1}(t)}{dt}+\frac{1}{C}\int[i_{1}%
(t)-i_{2}(t)]dt$

$\frac{1}{C}\int[i_{1}(t)-i_{2}(t)]dt=R_{2}i_{2}(t)$

$Vs(t)=R_{2}i_{2}(t)$

Transformada de Laplace:

$Ve(s)=R_{1}I_{1}(s)+LsI_{1}(s)+\frac{I_{1}(s)-I_{2}(s)}{Cs}$

$\frac{I_{1}(s)-I_{2}(s)}{Cs}=R_{2}I_{2}(s)$

$Vs(s)=R_{2}I_{2}(s)$

Despejar I$_{\text{1}}$(s):

$\frac{I_{1}(s)}{Cs}-\frac{I_{2}(s)}{Cs}=R_{2}I_{2}(s)$

$\frac{I_{1}(s)}{Cs}=R_{2}I_{2}(s)+\frac{I_{2}(s)}{Cs}$

$\frac{I_{1}(s)}{Cs}=(R_{2}+\frac{1}{Cs})I_{2}(s)$

$I_{1}(s)=(CR_{2}s+1)I_{2}(s)$

Sustituir I$_{\text{1}}$(s):

$Ve(s)=R_{1}(CR_{2}s+1)I_{2}(s)+Ls(CR_{2}s+1)I_{2}(s)+\frac{(CR_{2}%
s+1)I_{2}(s)-I_{2}(s)}{Cs}$

$Ve(s)=[CR_{1}R_{2}s+R_{1}+CR_{2}Ls^{2}+Ls+\frac{CR_{2}s+1}{Cs}-\frac{1}%
{Cs}]I_{2}$

$Ve=[CR_{2}Ls^{2}+(CR_{1}R_{2}+L)s+R_{1}+R_{2}]I_{2}$

Funci\'{o}n de transferencia:

$\frac{Vs(s)}{Ve(s)}=\frac{R_{2}I_{2}}{[CR_{2}Ls^{2}+(CR_{1}R_{2}%
+L)s+R_{1}+R_{2}]I_{2}}=\frac{R_{2}}{CR_{2}Ls^{2}+(CR_{1}R_{2}+L)s+R_{1}%
+R_{2}}$

Ecuaciones integrodiferenciales:

\newline

$R_{1}i_{1}(t)=Ve(t)-L\frac{di_{1}(t)}{dt}-\frac{1}{C}\int[i_{1}%
(t)-i_{2}(t)]dt$

$i_{1}(t)=[Ve(t)-L\frac{di_{1}(t)}{dt}-\frac{1}{C}\int[i_{1}(t)-i_{2}%
(t)]dt]\frac{1}{R_{1}}$

$i_{2}(t)=[\frac{1}{C}\int[i_{1}(t)-i_{2}(t)]dt]\frac{1}{R_{2}}$

Error estacionario:

$e(s)=\lim_{s\rightarrow0}sVe(s)[1-\frac{Ve(s)}{Ve(s)}]$

$e(s)=\lim_{s\rightarrow0}(s\times\frac{1}{s})[1-\frac{R_{2}}{CR_{2}%
Ls^{2}+(CR_{1}R_{2}+L)s+R_{1}+R_{2}}]$

$e(s)=1-\frac{R_{2}}{R_{1}+R_{2}}=\frac{R_{1}+R_{2}}{R_{1}+R_{2}}-\frac{R_{2}%
}{R_{1}+R_{2}}=\frac{R_{1}}{R_{1}+R_{2}}$

Control:

$\frac{10}{40+10}=\frac{1}{5}$

Caso:

$\frac{80}{10+80}=\frac{8}{9}$

Analisis de estabilidad de lazo abierto:

$\partial_{1,2}=\frac{-(CR_{1}R_{2}+L)\pm\sqrt{(CR_{1}R_{2}+L)^{2}%
-4(CR_{2}Ls)(R_{1}+R_{2})}}{2(CR_{2}Ls)}$

Control:

R1=10ohms, R2=40ohms, C=100uF, L=10uH

$\partial_{1}=-5.008\times10^{7}$

$\partial_{2}=-2.496\times10^{10}$

Caso:

R1=80ohms, R2=10ohms, C=20uF, L=10uF

$\partial_{1}=-7.9119\times10^{7}$

$\partial_{2}=-7.280\times10^{6}$

El sistema es estable dado que ambas raices son negativas y reales entonces el
sistema permite una respuesta sobreamortiguada
\end{thebibliography}


\end{document}